\documentclass[11pt,letterpaper]{article}
\usepackage[utf8]{inputenc}
\usepackage[T1]{fontenc}
\usepackage[french]{babel}
\usepackage{amsmath,amssymb,amsthm}
\usepackage{geometry}
\geometry{margin=2.5cm}
\usepackage{enumitem}
\usepackage{booktabs}
\usepackage{xcolor}
\usepackage{tcolorbox}
\usepackage{hyperref}

%% ------------------------------------------------------------
%% Environnements pour les notes de tableau
%% ------------------------------------------------------------
\newtheorem{definition}{Définition}[section]
\newtheorem{propriete}{Propriété}[section]
\newtheorem{remarque}{Remarque}[section]
\newtheorem{exemple}{Exemple}[section]

% Bloc de cours (~50 min)
\newcommand{\bloc}[1]{\subsection*{\fbox{Bloc #1}}}

% Encadré : ce qu'on écrit textuellement au tableau
\newtcolorbox{tableau}{colback=gray!8, colframe=black!60,
  title={\textbf{Au tableau}}, fonttitle=\bfseries}

% Encadré : moment R (portable/projecteur)
\newtcolorbox{momentR}{colback=blue!5, colframe=blue!50!black,
  title={\textbf{Moment R (projecteur)}}, fonttitle=\bfseries}

% Encadré : question à poser aux étudiants
\newcommand{\question}[1]{\medskip\noindent\textit{$\triangleright$ Question à la classe~: #1}\medskip}

% En-tête de semaine
\newcommand{\semaine}[2]{\newpage\section{Semaine #1 --- #2}}

\title{\textbf{STT-4300 --- Notes de tableau}\\[1ex]
\large Plan de cours hebdomadaire (14 semaines $\times$ 3 heures)\\
Rappels de régression linéaire, MLG, régression non-paramétrique et GAM}
\author{}
\date{}

\begin{document}
\maketitle

\noindent\textbf{Mode d'emploi.} Chaque semaine correspond à une séance de 3 heures,
découpée en trois blocs d'environ 50 minutes (avec pauses). Les encadrés
\textbf{Au tableau} contiennent ce qui est effectivement écrit au tableau~;
le texte entre les encadrés est le discours d'accompagnement (à dire, pas à écrire).
Les encadrés \textbf{Moment R} indiquent les démonstrations à faire au projecteur.

\tableofcontents

%% ============================================================
%% SEMAINE 1
%% ============================================================
\semaine{1}{Rappels de régression linéaire (I)~: modèle et estimation}

\noindent\textbf{Objectifs}~: (i) réactiver le modèle linéaire et son écriture
matricielle~; (ii) redériver l'estimateur des moindres carrés et ses propriétés~;
(iii) introduire la matrice chapeau et l'estimation de $\sigma^2$.

\bloc{1 --- Le modèle linéaire et son écriture matricielle}

Commencer par un exemple concret (p.~ex.\ \textsf{mpg} en fonction de la
puissance dans \textsf{Auto}) et demander aux étudiants de formuler le modèle.

\begin{tableau}
Données~: $\{(Y_i, X_i),\ i=1,\ldots,n\}$, $Y_i$ réponse continue,
$X_i = (X_{i1},\ldots,X_{ip})^\top$.

Modèle de régression linéaire multiple~:
$$Y_i = \beta_0 + \beta_1 X_{i1} + \cdots + \beta_p X_{ip} + \epsilon_i,
\qquad \epsilon_i \stackrel{iid}{\sim} \mathcal{N}(0,\sigma^2).$$

Écriture matricielle~: $Y = X\beta + \epsilon$, avec
$$X = \begin{pmatrix} 1 & X_{11} & \cdots & X_{1p} \\
\vdots & \vdots & & \vdots \\
1 & X_{n1} & \cdots & X_{np}\end{pmatrix}
\quad (n\times(p+1)), \qquad
\beta = (\beta_0,\beta_1,\ldots,\beta_p)^\top.$$

Conséquences des hypothèses~:
$$E[Y_i] = \mu_i = \beta_0 + \beta_1 X_{i1} + \cdots + \beta_p X_{ip},
\qquad Var[Y_i] = \sigma^2 \ \text{(homoscédasticité)}.$$
\end{tableau}

\question{quelles sont les deux hypothèses fortes de ce modèle~? (normalité
de la réponse~; linéarité de la relation). Annoncer que tout le cours consiste
à lever ces deux hypothèses, l'une après l'autre.}

\bloc{2 --- Estimation par moindres carrés (dérivation complète)}

\begin{tableau}
Critère des moindres carrés~:
$$l(\beta) = \sum_{i=1}^n \left\{Y_i - (\beta_0 + \beta_1 X_{i1} + \cdots
+ \beta_p X_{ip})\right\}^2 = (Y-X\beta)^\top(Y-X\beta).$$

Développer~:
$$l(\beta) = Y^\top Y - 2\beta^\top X^\top Y + \beta^\top X^\top X\beta.$$

Dériver par rapport à $\beta$ et annuler (équations normales)~:
$$\frac{\partial l}{\partial\beta} = -2X^\top Y + 2X^\top X\beta = 0
\quad\Longleftrightarrow\quad X^\top X\beta = X^\top Y.$$

Si $X^\top X$ inversible~:
$$\boxed{\hat\beta = (X^\top X)^{-1} X^\top Y}$$

Propriétés (à dériver au tableau)~:
\begin{align*}
E[\hat\beta] &= (X^\top X)^{-1}X^\top E[Y] = (X^\top X)^{-1}X^\top X\beta = \beta
\quad\text{(sans biais)}\\
Var[\hat\beta] &= (X^\top X)^{-1}X^\top\, \sigma^2 I\, X(X^\top X)^{-1}
= \sigma^2(X^\top X)^{-1}.
\end{align*}
Comme $Y$ gaussien et $\hat\beta$ linéaire en $Y$~:
$$\hat\beta \sim \mathcal{N}\left\{\beta,\ \sigma^2 (X^\top X)^{-1}\right\}.$$
\end{tableau}

Mentionner (sans démontrer)~: sous normalité, moindres carrés $=$ maximum de
vraisemblance~; théorème de Gauss--Markov (variance minimale parmi les
estimateurs linéaires sans biais).

\bloc{3 --- Matrice chapeau, résidus et estimation de $\sigma^2$}

\begin{tableau}
Valeurs prédites~:
$$\hat Y = X\hat\beta = \underbrace{X(X^\top X)^{-1}X^\top}_{=\,H}\,Y = HY.$$
$H$ = \textbf{matrice chapeau} (\textit{hat matrix})~; elle « met le chapeau »
sur $Y$. Propriétés~: $H^\top = H$, $H^2 = H$ (projection orthogonale sur
l'espace engendré par les colonnes de $X$).

Les éléments diagonaux $h_i = H_{ii}$ sont les \textbf{leviers}
(\textit{leverage})~; on y reviendra pour les diagnostics.

Résidus et estimation de $\sigma^2$~:
$$e_i = Y_i - \hat Y_i, \qquad SSE = \sum_{i=1}^n (Y_i - \hat Y_i)^2, \qquad
\hat\sigma^2 = \frac{SSE}{n-(p+1)}.$$
Le dénominateur $n-(p+1)$ = nombre d'observations $-$ nombre de paramètres
estimés (degrés de liberté).
\end{tableau}

\begin{momentR}
Ajuster une première régression avec \textsf{lm}, montrer \textsf{summary},
faire le lien entre chaque nombre affiché et les formules du tableau
(\textsf{Estimate}, \textsf{Std. Error}, \textsf{Residual standard error},
degrés de liberté).
\end{momentR}

%% ============================================================
%% SEMAINE 2
%% ============================================================
\semaine{2}{Rappels de régression linéaire (II)~: inférence, diagnostics et sélection}

\noindent\textbf{Objectifs}~: (i) inférence sur les coefficients (tests $t$, test $F$)~;
(ii) prédiction avec intervalles~; (iii) diagnostics graphiques et mesures d'influence~;
(iv) critères AIC/BIC~; (v) motiver la suite du cours (limites du modèle linéaire).

\bloc{1 --- Inférence statistique}

\begin{tableau}
IC à $100(1-\alpha)\%$ pour $\beta_j$~:
$$\left[\hat\beta_j \pm t_{\alpha/2,\,n-(p+1)}\, se(\hat\beta_j)\right],
\qquad se(\hat\beta_j) = \sqrt{\hat\sigma^2 \left[(X^\top X)^{-1}\right]_{j+1,j+1}}.$$

Test de $H_0: \beta_j = 0$ vs $H_1: \beta_j\neq 0$~:
$$t_{obs} = \frac{\hat\beta_j}{se(\hat\beta_j)} \sim t_{n-(p+1)} \text{ sous } H_0~;
\quad\text{rejet si } |t_{obs}| > t_{\alpha/2,\,n-(p+1)}.$$

\textbf{Modèles emboîtés}~: le modèle réduit ($q+1$ param.) est un cas particulier
du modèle complet ($p+1$ param., $q<p$). Test $F$~:
$$F_{obs} = \frac{(SSE_{r\acute{e}duit} - SSE_{complet})/(p-q)}
{SSE_{complet}/(n-(p+1))} \ \sim\ F_{p-q,\ n-(p+1)} \text{ sous } H_0.$$

Coefficient de détermination~:
$$R^2 = 1 - \frac{SSE}{SST}, \qquad
R^2_{adj} = 1 - \frac{SSE/(n-(p+1))}{SST/(n-1)}, \qquad
SST = \sum_i(Y_i-\bar Y)^2.$$
\end{tableau}

\question{pourquoi le $R^2$ augmente-t-il toujours quand on ajoute une variable~?
D'où la nécessité du $R^2$ ajusté et, plus tard, des critères pénalisés.}

\bloc{2 --- Prédiction et diagnostics}

\begin{tableau}
Nouvelle observation $X^* = (1, X_1^*,\ldots,X_p^*)^\top$~:
$\hat Y^* = {X^*}^\top\hat\beta$, $Var[\hat Y^*] = \sigma^2\,{X^*}^\top(X^\top X)^{-1}X^*$.

Deux intervalles \textbf{à ne pas confondre}~:
\begin{itemize}
\item IC pour la \textbf{moyenne} $E[Y^*]$~:
$\ {X^*}^\top\hat\beta \pm t_{\alpha/2,\,n-(p+1)}\,\hat\sigma
\sqrt{{X^*}^\top (X^\top X)^{-1} X^*}$~;
\item Intervalle de \textbf{prédiction} pour $Y^*$~:
$\ {X^*}^\top\hat\beta \pm t_{\alpha/2,\,n-(p+1)}\,\hat\sigma
\sqrt{1 + {X^*}^\top (X^\top X)^{-1} X^*}$.
\end{itemize}
Le « $1+$ » vient de la variance de l'erreur $\epsilon^*$ de la nouvelle observation.

\textbf{Résidus}~: bruts $e_i = Y_i - \hat Y_i$~; standardisés
$r_i = e_i/(\hat\sigma\sqrt{1-h_i})$~; studentisés ($\sigma$ estimé sans
l'observation $i$).

Vérifications graphiques (4 graphiques à esquisser au tableau)~:
linéarité (résidus vs $\hat Y$), homoscédasticité, normalité (QQ-plot),
indépendance (résidus vs ordre).

\textbf{Influence}~: levier $h_i$ (seuil indicatif $2(p+1)/n$)~; distance de Cook
$$D_i = \frac{r_i^2}{p+1}\cdot\frac{h_i}{1-h_i}~;$$
DFBETAS (effet du retrait de $i$ sur chaque $\hat\beta_j$),
DFFITS (effet sur $\hat Y_i$).
\end{tableau}

\begin{momentR}
Montrer \textsf{plot(modele)} (les 4 graphiques de diagnostic),
\textsf{predict(..., interval="confidence")} vs
\textsf{interval="prediction"}, \textsf{cooks.distance}.
\end{momentR}

\bloc{3 --- Sélection de variables et limites du modèle linéaire}

\begin{tableau}
Critères d'information (valables pour des modèles emboîtés ou non)~:
$$AIC = -2\, l(\hat\beta;\, y) + 2\,(p+1), \qquad
BIC = -2\, l(\hat\beta;\, y) + \log(n)\,(p+1).$$
Plus petit $=$ meilleur. Le BIC pénalise plus fortement dès que $n\geq 8$.

Procédures automatiques~: \textit{forward} (ajout progressif),
\textit{backward} (retrait progressif), \textit{stepwise} (mixte).
En \textrm{R}~: \textsf{step}.
\end{tableau}

Conclure la semaine en écrivant les \textbf{deux limites} du modèle linéaire~:
\begin{tableau}
Le modèle linéaire suppose~:
\begin{enumerate}
\item $Y$ continue et normale \quad $\to$ que faire si $Y$ est binaire~?
un dénombrement~? \quad (chapitres 2--3)~;
\item $E[Y]$ linéaire en $X$ \quad $\to$ que faire si la relation est
fortement non linéaire~? \quad (chapitres 6--8).
\end{enumerate}
\end{tableau}

\question{donnez-moi des exemples de réponses non normales dans vos domaines
(défaut de paiement, présence d'une maladie, nombre de réclamations\ldots).
Noter les exemples au tableau~: ils serviront toute la session.}

%% ============================================================
%% SEMAINE 3
%% ============================================================
\semaine{3}{Données binaires (I)~: modèle, fonctions de lien et estimation}

\noindent\textbf{Objectifs}~: (i) présenter les données Bernoulli et binomiales par
des exemples~; (ii) introduire prédicteur linéaire et fonction de lien~;
(iii) interpréter les coefficients (rapport de cotes)~; (iv) écrire la
log-vraisemblance et expliquer l'estimation numérique.

\bloc{1 --- Exemples et structure des données}

Présenter les quatre exemples de \texttt{GLMsData}~: \textsf{quilpie} (pluie $>$ 10 mm~:
oui/non), \textsf{nminer} (oiseau détecté~: oui/non), \textsf{deposit} (insectes tués
parmi $m_i$ exposés), \textsf{turbines} (machines fissurées parmi $m_i$).

\begin{tableau}
\textbf{Données individuelles (Bernoulli)}~: $Y_i\in\{0,1\}$,
$$P(Y_i = y_i) = \pi_i^{y_i}(1-\pi_i)^{1-y_i}, \qquad
E[Y_i] = \pi_i, \qquad Var[Y_i] = \pi_i(1-\pi_i).$$

\textbf{Données groupées (binomiale)}~: $Y_i\in\{0,\ldots,m_i\}$, $m_i$ connu,
$$P(Y_i = y_i) = \binom{m_i}{y_i}\pi_i^{y_i}(1-\pi_i)^{m_i-y_i}, \qquad
E[Y_i] = m_i\pi_i, \qquad Var[Y_i] = m_i\pi_i(1-\pi_i).$$

Bernoulli $=$ binomiale avec $m_i = 1$~: un seul cadre pour les deux.

\textbf{Remarque clé}~: la variance dépend de la moyenne~!
(contrairement au modèle linéaire).
\end{tableau}

\question{pourquoi ne pas simplement faire une régression linéaire de $Y_i$
sur les $X$~? Laisser chercher~: prédictions hors de $[0,1]$, variance non
constante, normalité impossible.}

\bloc{2 --- Prédicteur linéaire et fonctions de lien}

\begin{tableau}
Prédicteur linéaire~: $\eta_i = \beta_0 + \beta_1 X_{i1} + \cdots + \beta_p X_{ip}$.

Fonction de lien $g$ (connue, monotone)~: $\eta_i = g(\pi_i)$, i.e.
$\pi_i = g^{-1}(\eta_i) \in [0,1]$.

Trois choix classiques~:
\begin{itemize}
\item \textbf{logit}~: $g(u) = \log\{u/(1-u)\}$
\ $\Rightarrow$ \ $\pi_i = \dfrac{e^{\eta_i}}{1+e^{\eta_i}}$~;
\item \textbf{probit}~: $g(u) = \Phi^{-1}(u)$
\ $\Rightarrow$ \ $\pi_i = \Phi(\eta_i)$~;
\item \textbf{c-log-log}~: $g(u) = \log\{-\log(1-u)\}$
\ $\Rightarrow$ \ $\pi_i = 1 - e^{-e^{\eta_i}}$.
\end{itemize}
Dans les trois cas, $g^{-1}$ est une fonction de répartition
(tracer les trois courbes en S au tableau~: très similaires en pratique).

\textbf{Interprétation (lien logit)}~: cote (\textit{odds})
$o_i = \pi_i/(1-\pi_i) = e^{\eta_i}$.
$X_j \to X_j + 1$ \ $\Rightarrow$ \ cote multipliée par $e^{\beta_j}$
= \textbf{rapport de cotes} (\textit{odds ratio}).
\end{tableau}

Faire un exemple numérique d'interprétation~: si $\hat\beta_1 = 0.7$, alors
$e^{0.7}\approx 2$~: la cote double pour chaque unité de $X_1$.

\bloc{3 --- Estimation par maximum de vraisemblance}

\begin{tableau}
Log-vraisemblance (cas binomial, couvre Bernoulli)~:
$$l = \sum_{i=1}^n\left\{\log\binom{m_i}{y_i} + y_i\log(\pi_i)
+ (m_i-y_i)\log(1-\pi_i)\right\},$$
avec $\pi_i = g^{-1}(\beta_0 + \beta_1 X_{i1} + \cdots + \beta_p X_{ip})$.

En posant $\mu_i = m_i\pi_i$~:
$$l(\mu;y) = \sum_{i=1}^n\left\{y_i\log(\mu_i) + (m_i-y_i)\log(m_i-\mu_i)\right\} + K.$$

Pas de forme fermée pour $\hat\beta$ \ $\to$ \ maximisation numérique
(moindres carrés pondérés itératifs~; \textsf{glm} en \textrm{R}).

Matrice d'information de Fisher $I$~: $I_{jk} = -\partial^2 l/\partial\beta_j\partial\beta_k$
évaluée en $\hat\beta$. Alors
$$v(\hat\beta) = I^{-1} = (X^\top WX)^{-1}, \qquad
W = \mathrm{diag}\{w_i\},\quad w_i = \frac{\hat\mu_i(m_i-\hat\mu_i)}{m_i},$$
et asymptotiquement $\hat\beta \approx \mathcal{N}\{\beta,\ v(\hat\beta)\}$.
\end{tableau}

\begin{momentR}
Premier \textsf{glm}~: \textsf{glm(y \textasciitilde{} SOI, family=binomial,
data=quilpie)}. Comparer la sortie \textsf{summary} à celle de \textsf{lm}~:
\textsf{z value} au lieu de \textsf{t value}, déviance au lieu de $SSE$
(annoncer que tout cela sera justifié aux semaines 4 et 6).
Montrer la syntaxe \textsf{cbind(succès, échecs)} pour les données groupées.
\end{momentR}

%% ============================================================
%% SEMAINE 4
%% ============================================================
\semaine{4}{Données binaires (II)~: inférence, déviance et sélection de modèles}

\noindent\textbf{Objectifs}~: (i) tests et IC de Wald~; (ii) IC pour une probabilité
prédite (deux méthodes)~; (iii) déviance et comparaison de modèles emboîtés~;
(iv) AIC/BIC et sélection automatique.

\bloc{1 --- Inférence de Wald}

\begin{tableau}
IC à $100(1-\alpha)\%$ pour $\beta_j$ (Wald, asymptotique)~:
$$\left[\hat\beta_j \pm z_{\alpha/2}\sqrt{v(\hat\beta_j)}\right],
\qquad z_\gamma~: P(Z>z_\gamma)=\gamma,\ Z\sim\mathcal{N}(0,1).$$

Test de $H_0:\beta_j = 0$~: $z_{obs} = \hat\beta_j/\sqrt{v(\hat\beta_j)}$~;
rejet si $|z_{obs}| > z_{\alpha/2}$~;
$p$-valeur $= 2P(Z > |z_{obs}|)$.

Différence avec la régression linéaire~: loi normale (asymptotique) au lieu de
Student (exacte).
\end{tableau}

\bloc{2 --- IC pour une probabilité prédite~; exemple \textsf{turbines}}

\begin{tableau}
Pour $X^*$ donné~: $\hat\pi^* = g^{-1}({X^*}^\top\hat\beta)$.
Deux méthodes pour un IC~:

\textbf{(1) Transformation}~: IC pour $\eta^* = {X^*}^\top\beta$~:
$$\left[{X^*}^\top\hat\beta \pm z_{\alpha/2}\sqrt{{X^*}^\top v(\hat\beta)X^*}\right],$$
puis appliquer $g^{-1}$ aux deux bornes \ $\to$ \ toujours dans $[0,1]$.

\textbf{(2) Règle du delta}~:
$v(\hat\pi^*) = h({X^*}^\top\hat\beta)\, v({X^*}^\top\hat\beta)$ avec
$h(u) = \{\partial g^{-1}(u)/\partial u\}^2$, puis
$\left[\hat\pi^* \pm z_{\alpha/2}\sqrt{v(\hat\pi^*)}\right]$.

En pratique~: résultats très similaires. Préférer (1) près de 0 ou 1.
\end{tableau}

\begin{momentR}
Exemple \textsf{turbines}~: ajuster les trois liens (logit, probit, cloglog),
superposer les trois courbes $\hat\pi^*$ (quasi identiques), calculer les IC
avec \textsf{predict(..., se.fit=TRUE)} sur l'échelle du lien puis transformer.
\end{momentR}

\bloc{3 --- Déviance et sélection de modèles}

\begin{tableau}
\textbf{Modèle saturé}~: $n$ paramètres $\tilde\mu_i$, EMV $= y_i$
(ajustement parfait). \textbf{Déviance} du modèle considéré~:
$$D = 2\{l(y,y) - l(\hat\mu,y)\}
= 2\sum_{i=1}^n\left\{y_i\log\frac{y_i}{\hat\mu_i}
+ (m_i-y_i)\log\frac{m_i-y_i}{m_i-\hat\mu_i}\right\} = \sum_i d_i.$$
$D$ joue le rôle du $SSE$~: petit $D$ = bon ajustement.

\textbf{Modèles emboîtés}~: $D_2$ (réduit) $-$ $D_1$ (complet)~:
$$\chi^2 = D_2 - D_1 \sim \chi^2_d \text{ sous } H_0
\quad (d = \text{diff.\ du nombre de paramètres}).$$

\textbf{Modèles non emboîtés}~:
$$AIC = -2\,l(\hat\mu;y) + 2(p+1), \qquad
BIC = -2\,l(\hat\mu;y) + \log(n)(p+1).$$
Sélection automatique \textit{forward}/\textit{backward}/\textit{stepwise}~:
comme en régression linéaire (\textsf{step}).
\end{tableau}

\begin{momentR}
Sur \textsf{turbines}~: tester $H_0:\beta_2=\beta_3=0$ (termes en
\textsf{Hours}$^2$, \textsf{Hours}$^3$) par différence de déviances~:
\textsf{1 - pchisq(stat, df=2)} donne $0.4767$ $\to$ modèle réduit adéquat.
Sur \textsf{deposit}~: \textsf{step} retient type d'insecticide $+$ dose $+$ dose$^2$.
\end{momentR}

%% ============================================================
%% SEMAINE 5
%% ============================================================
\semaine{5}{Données binaires (III)~: diagnostics, surdispersion et courbe ROC}

\noindent\textbf{Objectifs}~: (i) résidus déviance et de Pearson~; (ii) tests
d'adéquation globale et leurs limites~; (iii) surdispersion et quasi-binomiale~;
(iv) sensibilité, spécificité, courbe ROC et AUC.

\bloc{1 --- Résidus et adéquation}

\begin{tableau}
\textbf{Résidus déviance}~: $r_{D,i} = \mathrm{signe}(y_i-\hat\mu_i)\sqrt{d_i}$.

\textbf{Résidus de Pearson}~:
$$r_{P,i} = \frac{y_i - \hat\mu_i}{\sqrt{\hat\mu_i(m_i-\hat\mu_i)/m_i}}
\quad\text{(observation centrée-réduite)}.$$

\textbf{Adéquation globale}~: si le modèle est bien spécifié et
$m_i\to\infty$~:
$$D = \sum_i r_{D,i}^2 \quad\text{et}\quad
\chi^2_P = \sum_i r_{P,i}^2 \ \approx\ \chi^2_{n-(p+1)}.$$

\textbf{Attention}~: valide seulement pour données \textbf{groupées} avec
$m_i$ élevés. Inutilisable pour données Bernoulli~!
\end{tableau}

\question{pourquoi le test d'adéquation échoue-t-il avec des données 0/1~?
(l'asymptotique est en $m_i$, pas en $n$).}

\bloc{2 --- Extra-variabilité (surdispersion)}

\begin{tableau}
Diagnostic~: plusieurs $|r_i| > 3$ \ $\to$ \ suspicion de
\textbf{surdispersion}~: $Var[Y_i] > m_i\pi_i(1-\pi_i)$.

Modélisation~: $Var[Y_i] = \phi\, m_i\pi_i(1-\pi_i)$, avec
$$\hat\phi = \frac{\chi^2_P}{n-(p+1)}.$$
En \textrm{R}~: \textsf{family=quasibinomial()} — mêmes $\hat\beta$,
erreurs types multipliées par $\sqrt{\hat\phi}$.

Vérification de la linéarité du prédicteur~: tracer
$\{(\hat\pi_i, \tilde\pi_i = y_i/m_i)\}$~: points autour de $y=x$ si OK.

Observations influentes~: leviers, Cook, DFBETAS, DFFITS — mêmes
interprétations qu'en régression linéaire.
\end{tableau}

\bloc{3 --- Spécificité, sensibilité et courbe ROC}

Motiver~: pour des données individuelles, l'adéquation est difficile à évaluer,
mais on peut mesurer la \textbf{capacité prédictive}.

\begin{tableau}
Seuil $u\in[0,1]$ (souvent $1/2$)~: classer $\hat y_i = \mathbf{1}_{\{\hat\pi_i > u\}}$.

$$\text{Spécificité} = P(\hat y = 0\mid y = 0), \qquad
\text{Sensibilité} = P(\hat y = 1\mid y = 1).$$

Estimateurs (à écrire avec les indicatrices)~:
$$\widehat{Spec}(u) = \frac{\#\{\hat\pi_i < u \text{ et } y_i = 0\}}{\#\{y_i=0\}},
\qquad
\widehat{Sens}(u) = \frac{\#\{\hat\pi_i > u \text{ et } y_i = 1\}}{\#\{y_i=1\}}.$$

\textbf{Courbe ROC}~: faire varier $u$ de 0 à 1, tracer
$\{(1-\widehat{Spec}(u),\ \widehat{Sens}(u))\}$.
Diagonale $=$ classement au hasard.

\textbf{AUC} (aire sous la courbe)~:
$[0.9,1]$ excellent~; $[0.8,0.9)$ bon~; $[0.7,0.8)$ moyen~;
$[0.6,0.7)$ mauvais~; $[0.5,0.6)$ très faible.
\end{tableau}

\begin{momentR}
Sur \textsf{quilpie}~: \textsf{roc.curve} (librairie \texttt{clinfun}) avec les
probabilités ajustées~; AUC $= 0.784$ (e.t.\ $0.057$) $\to$ bon modèle.
Faire varier le seuil $u$ en direct pour montrer le compromis
sensibilité/spécificité.
\end{momentR}

%% ============================================================
%% SEMAINE 6
%% ============================================================
\semaine{6}{Modèles linéaires généralisés (I)~: la famille exponentielle}

\noindent\textbf{Objectifs}~: (i) définir la famille exponentielle~; (ii) dériver
espérance et variance à partir de $b(\theta)$~; (iii) reconnaître normale,
binomiale, Poisson et gamma comme cas particuliers~; (iv) définir la déviance
dans le cadre général.

\bloc{1 --- Définition et propriétés fondamentales}

Motiver~: régressions linéaire et logistique partagent la même mécanique
(coefficients, tests, sélection, diagnostics) $\to$ chercher le cadre unifié.

\begin{tableau}
\textbf{Définition.} $Y$ appartient à la \textbf{famille exponentielle} si
$$f_Y(y) = a(y,\phi)\exp\left\{\frac{y\theta - b(\theta)}{\phi}\right\},$$
$\theta$ = paramètre \textbf{canonique}, $\phi$ = paramètre de \textbf{dispersion}.

\textbf{Propriété.} $E[Y] = \mu = b'(\theta)$ et
$Var[Y] = \phi\, b''(\theta) = \phi\, V(\mu)$,
où $V(u) = b''[(b')^{-1}(u)]$ est la \textbf{fonction de variance}.

Esquisse de preuve (à faire)~: partir de $\int f_Y(y)\,dy = 1$, dériver sous
l'intégrale par rapport à $\theta$~:
$$\int \frac{y - b'(\theta)}{\phi}f_Y(y)\,dy = 0 \ \Rightarrow\ E[Y] = b'(\theta)~;$$
dériver une seconde fois pour la variance.
\end{tableau}

\bloc{2 --- Exemples (à dériver un par un au tableau)}

\begin{tableau}
\textbf{Normale} $\mathcal{N}(\mu,\sigma^2)$~: réécrire la densité~:
$$f_Y(y) = \underbrace{\tfrac{1}{\sqrt{2\pi\sigma^2}}e^{-y^2/(2\sigma^2)}}_{a(y,\phi)}
\exp\left\{\frac{y\mu - \mu^2/2}{\sigma^2}\right\}
\ \Rightarrow\ \theta=\mu,\ \phi=\sigma^2,\ b(\theta)=\theta^2/2,\ V(\mu)=1.$$

\textbf{Binomiale} $\mathrm{Bin}(m,\pi)$~:
$\theta = \log\frac{\pi}{1-\pi}$, $\phi=1$, $b(\theta) = m\log(1+e^\theta)$,
$V(\mu) = \mu(m-\mu)/m$.

\textbf{Poisson} $\mathcal{P}(\lambda)$~:
$\theta = \log\lambda$, $\phi = 1$, $b(\theta) = e^\theta$, $V(\mu) = \mu$.

\textbf{Gamma} (reparamétrée, $\mu = \alpha\beta$, $\phi = 1/\alpha$)~:
$\theta = -1/\mu$, $b(\theta) = -\log(-\theta)$, $V(\mu) = \mu^2$.

\medskip
\begin{tabular}{lcccc}
\toprule
Distribution & $\theta$ & $\phi$ & $b(\theta)$ & $V(\mu)$ \\
\midrule
$\mathcal{N}(\mu,\sigma^2)$ & $\mu$ & $\sigma^2$ & $\theta^2/2$ & $1$ \\
$\mathrm{Bin}(m,\pi)$ & $\log\{\pi/(1-\pi)\}$ & $1$ & $m\log(1+e^\theta)$ & $\mu(m-\mu)/m$ \\
$\mathcal{P}(\lambda)$ & $\log(\lambda)$ & $1$ & $e^\theta$ & $\mu$ \\
$\mathrm{Gamma}(\mu,\phi)$ & $-1/\mu$ & $\phi$ & $-\log(-\theta)$ & $\mu^2$ \\
\bottomrule
\end{tabular}
\end{tableau}

Faire vérifier par les étudiants (exercice de 10 min en classe)~: retrouver
$E[Y]$ et $Var[Y]$ de la Poisson à partir de $b(\theta) = e^\theta$.

\bloc{3 --- Déviance dans le cadre général}

\begin{tableau}
Poser $t(y,\mu) = y\theta - b(\theta)$ et
$d(y,\mu) = 2\{t(y,y) - t(y,\mu)\}$.

Montrer que $d(y,\mu)\geq 0$ avec égalité ssi $\mu = y$~:
$$\frac{\partial d}{\partial\theta} = -\{y - b'(\theta)\}, \qquad
\frac{\partial^2 d}{\partial\theta^2} = b''(\theta) = V(\mu) > 0.$$

\textbf{Cas normal}~: $d(y,\mu) = (y-\mu)^2$
$\to$ la déviance \textbf{généralise le carré du résidu}.

Pour $Y_1,\ldots,Y_n$ indépendantes (dispersion commune $\phi$)~:
$$D(y,\mu) = \sum_{i=1}^n d(y_i,\mu_i), \qquad
D^*(y,\mu) = \frac{D(y,\mu)}{\phi} \stackrel{approx.}{\sim} \chi^2_n.$$
\end{tableau}

Terminer en annonçant la définition complète du MLG (semaine 7)~: les trois
composantes seront la famille exponentielle, le prédicteur linéaire, le lien.

%% ============================================================
%% SEMAINE 7
%% ============================================================
\semaine{7}{Modèles linéaires généralisés (II)~: définition, estimation et inférence}

\noindent\textbf{Objectifs}~: (i) définir formellement un MLG (trois composantes)~;
(ii) estimation de $\beta$ et de $\phi$~; (iii) inférence pour $\phi$ connu et
inconnu~; (iv) sélection et validation dans le cadre général.

\bloc{1 --- Les trois composantes d'un MLG}

\begin{tableau}
Un \textbf{modèle linéaire généralisé} $=$ trois composantes~:
\begin{enumerate}
\item \textbf{Aléatoire}~: $Y_i \sim$ famille exponentielle$(\theta_i,\phi)$
($\phi$ commun)~: $E[Y_i] = \mu_i = b'(\theta_i)$,
$Var[Y_i] = \phi V(\mu_i)$~;
\item \textbf{Systématique}~:
$\eta_i = \beta_0 + \beta_1 X_{i1} + \cdots + \beta_p X_{ip}$~;
\item \textbf{Lien}~: $g(\mu_i) = \eta_i$, $g$ monotone dérivable.
Lien \textbf{canonique} si $\eta_i = \theta_i$, i.e.\ $g[b'(u)] = u$.
\end{enumerate}

Cas particuliers (tableau à compléter avec la classe)~:
\begin{tabular}{lll}
\toprule
Modèle & Distribution & Lien canonique \\
\midrule
Régression linéaire & normale & identité \\
Régression logistique & binomiale & logit \\
Régression Poisson & Poisson & log \\
\bottomrule
\end{tabular}
\end{tableau}

\bloc{2 --- Estimation de $\beta$ et de $\phi$}

\begin{tableau}
Log-vraisemblance~:
$$l(\beta,\phi\mid\Delta)
= \underbrace{\sum_i\log a(y_i,\phi)}_{l_1(\phi)}
+ \frac{1}{\phi}\underbrace{\sum_i\{y_i\theta_i - b(\theta_i)\}}_{l_2(\beta)}.$$
$\hat\beta$ maximise $l_2(\beta)$~: il \textbf{ne dépend pas de $\phi$}.

Distribution asymptotique~:
$$\hat\beta \sim \mathcal{N}\left\{\beta,\ \phi(X^\top WX)^{-1}\right\},
\qquad W_i = \frac{1}{V(\hat\mu_i)[g'(\hat\mu_i)]^2}.$$

\textbf{Estimation de $\phi$} (si nécessaire~: normale, gamma)~:
par la méthode des moments, à partir de
$$X^2 = \sum_i\frac{(y_i-\hat\mu_i)^2}{V(\hat\mu_i)}
\quad\text{ou}\quad D(y,\hat\mu),$$
qui, divisées par $\phi$, sont $\approx\chi^2_{n-(p+1)}$~:
$$\tilde\phi_P = \frac{X^2}{n-(p+1)} \ \text{(Pearson, défaut de \textsf{glm})},
\qquad \tilde\phi_D = \frac{D(y,\hat\mu)}{n-(p+1)}.$$
Pearson~: biais négligeable, robuste. Déviance~: faible variance mais sensible
aux $y$ proches des bornes.
\end{tableau}

\bloc{3 --- Inférence, sélection et validation}

\begin{tableau}
\textbf{Cas $\phi$ connu} (logistique, Poisson)~: quantiles normaux~:
$$\text{IC~: } \left[\hat\beta_j \pm z_{\alpha/2}\sqrt{\phi\, v_j}\right],
\quad v_j = [(X^\top WX)^{-1}]_{j+1,j+1}~;$$
modèles emboîtés~:
$X_{obs} = \{D_A - D_B\}/\phi \sim \chi^2_d$ sous $H_0$.

\textbf{Cas $\phi$ inconnu}~: remplacer $\phi$ par $\tilde\phi_P$ et
normale par Student~:
$$\left[\hat\beta_j \pm t_{\alpha/2,\,n-(p+1)}\sqrt{\tilde\phi_P\, v_j}\right],
\qquad
X_{obs} = \frac{(D_A - D_B)/d}{\tilde\phi_P} \sim F_{d,\,n-(p+1)}.$$
(Parallèle exact avec $z\to t$ et $\chi^2 \to F$ en régression linéaire.)

\textbf{Sélection}~: AIC/BIC $+$ \textit{forward}/\textit{backward}/%
\textit{stepwise}, comme avant.
\textbf{Validation}~: résidus, leviers, Cook, DFBETAS, DFFITS~: tout reste valide.
\end{tableau}

\begin{momentR}
Reprendre les modèles des semaines 1 et 3 et montrer qu'ils sont deux instances
de \textsf{glm} (\textsf{family=gaussian} et \textsf{family=binomial})~:
mêmes sorties, même mécanique. Ajuster un modèle gamma
(\textsf{family=Gamma}) pour montrer un cas où $\phi$ est estimé
(ligne \textsf{Dispersion parameter} de \textsf{summary}).
\end{momentR}

%% ============================================================
%% SEMAINE 8
%% ============================================================
\semaine{8}{MLG (III)~: régression Poisson, offset et surdispersion}

\noindent\textbf{Objectifs}~: (i) régression Poisson et interprétation multiplicative~;
(ii) variables offset pour modéliser des taux~; (iii) surdispersion~:
quasi-Poisson et binomiale négative.

\bloc{1 --- La régression Poisson}

\begin{tableau}
$Y_i\in\mathbb{N}$ (dénombrement)~: $Y_i\sim\mathcal{P}(\mu_i)$,
$$P(Y_i=y_i) = e^{-\mu_i}\frac{\mu_i^{y_i}}{y_i!}, \qquad
E[Y_i] = Var[Y_i] = \mu_i.$$

Lien canonique $g(\mu) = \log(\mu)$~:
$$\mu_i = e^{\beta_0}\times\{e^{\beta_1}\}^{X_{i1}}\times\cdots\times
\{e^{\beta_p}\}^{X_{ip}}.$$
\textbf{Interprétation multiplicative}~: $X_j \to X_j+1$
$\Rightarrow$ $\mu_i$ multiplié par $e^{\beta_j}$.

Log-vraisemblance~:
$$l(\beta) = \sum_i\{y_i\,\eta_i - e^{\eta_i}\} + K, \qquad
\hat\beta\approx\mathcal{N}\{\beta, (X^\top WX)^{-1}\},\ w_i = \hat\mu_i.$$

Déviance individuelle et adéquation~:
$$d(y,\mu) = 2\left\{y\log\frac{y}{\mu} - (y-\mu)\right\}, \qquad
X^2 = \sum_i\frac{(y_i-\hat\mu_i)^2}{\hat\mu_i}
\ \approx\ \chi^2_{n-(p+1)}.$$
\end{tableau}

\bloc{2 --- Variables offset~: modéliser des taux}

Présenter \textsf{danishlc}~: cas de cancer du poumon par ville (4) et tranche
d'âge (6), population $m_i$ connue par strate.

\question{pourquoi pas une régression logistique ici~?
($\tilde\pi \approx 3.6\times10^{-3}$~: événement rare, et on connaît surtout
la population exposée, pas des essais individuels).}

\begin{tableau}
On modélise le \textbf{taux} $\mu_i/m_i$~:
$$\log\left(\frac{\mu_i}{m_i}\right) = \beta_0 + \beta_1 X_{i1}+\cdots+\beta_p X_{ip}
\ \Longleftrightarrow\
\log(\mu_i) = \underbrace{\log(m_i)}_{\text{offset}} + \beta_0 + \cdots + \beta_pX_{ip}.$$
$\log(m_i)$~: variable explicative à coefficient \textbf{fixé à 1} = \textbf{offset}.

Utile dès que $E[Y_i]\propto$ exposition (population, durée, effort).
\end{tableau}

\begin{momentR}
\textsf{glm(Cases \textasciitilde{} offset(log(Pop)) + City + Age,
family=poisson, data=danishlc)}. Tests par différence de déviances~:
effet d'âge significatif, pas d'effet de ville.
\end{momentR}

\bloc{3 --- Variabilité extra-Poissonienne}

\begin{tableau}
Problème fréquent~: $Var[Y_i] > E[Y_i]$ (\textbf{surdispersion}).

\textbf{(1) Quasi-Poisson}~: on ne spécifie que les moments~:
$E[Y_i] = \mu_i$, $Var[Y_i] = \phi\mu_i$.
Mêmes $\hat\beta$, variance multipliée par $\tilde\phi_P$.
\textbf{Pas de vraisemblance} $\Rightarrow$ pas d'AIC/BIC.
En \textrm{R}~: \textsf{family=quasipoisson}.

\textbf{(2) Binomiale négative}~: $E[Y_i] = \mu_i$,
$Var[Y_i] = \mu_i + \alpha\mu_i^2$, $\alpha>0$~;
Poisson retrouvée quand $\alpha\to0$.
Interprétation des $\beta$ inchangée. En \textrm{R}~: \textsf{glm.nb} (\texttt{MASS}).

Test de $H_0: \alpha = 0$ (frontière du domaine~!)~:
$Q_{obs} = D_0 - D_1$ et
$$Q \sim \tfrac12\chi^2_0 + \tfrac12\chi^2_1 \text{ sous } H_0,
\qquad p\text{-valeur} = \tfrac12 P(\chi^2_1 > Q_{obs}).$$
\end{tableau}

\begin{momentR}
Sur \textsf{danishlc}~: comparer Poisson et binomiale négative~;
$p$-valeur $= 0.143$ $\to$ pas de surdispersion détectée.
\end{momentR}

Conclure la première moitié du cours~: on sait maintenant traiter des réponses
normales, binaires, de dénombrement et gamma \textbf{avec un prédicteur
linéaire}. Reste la deuxième limite~: la \textbf{linéarité}~; ce sera l'objet
des semaines 9 à 14, en commençant par la question « comment comparer des
modèles selon leur capacité de \textbf{prédiction}~? ».

%% ============================================================
%% SEMAINE 9
%% ============================================================
\semaine{9}{Prédiction, taux d'erreur et validation croisée}

\noindent\textbf{Objectifs}~: (i) distinguer adéquation et précision de prédiction~;
(ii) approche par ensemble de validation~; (iii) LOOCV et sa formule explicite~;
(iv) validation croisée à $K$ blocs~; (v) mesures de performance en classification.

\bloc{1 --- Adéquation versus prédiction}

Présenter \textsf{Auto} ($n=392$)~: \texttt{mpg} vs \textsf{horsepower},
et les quatre modèles candidats.

\begin{tableau}
Modèles candidats~:
\begin{align*}
\text{M1~: } Y_i &= \beta_0 + \beta_1 X_i + \epsilon_i\\
\text{M2~: } Y_i &= \beta_0 + \beta_1 X_i + \beta_2\log(X_i) + \epsilon_i\\
\text{M3~: } Y_i &= \beta_0 + \beta_1 X_i + \beta_2 X_i^2 + \epsilon_i\\
\text{M4~: } Y_i &= \beta_0 + \beta_1 X_i + \beta_2 X_i^2 + \beta_3 X_i^3 + \epsilon_i
\end{align*}

Deux critères \textbf{différents}~:
\begin{itemize}
\item \textbf{Adéquation} aux données observées~: $SSE$, AIC, BIC, tests
$\to$ \textbf{taux d'erreur d'entraînement}~;
\item \textbf{Prédiction} d'une nouvelle observation $(X_0,Y_0)$~:
minimiser $E[(Y_0 - \hat Y_0)^2]$
$\to$ \textbf{taux d'erreur test}.
\end{itemize}

Cas extrême~: avec $n$ points, un polynôme de degré $n-1$ donne $SSE = 0$
mais prédit très mal~: \textbf{surajustement}.
\end{tableau}

Dessiner au tableau la courbe classique~: erreur d'entraînement qui décroît
toujours avec la complexité~; erreur test en U.

\bloc{2 --- Validation croisée}

\begin{tableau}
\textbf{(1) Ensemble de validation}~: partager les données en entraînement
($n_1$ obs.)\ et test ($n_2 = n - n_1$)~:
$$\widehat{Err}_{test} = \frac{1}{n_2}\sum_{i=n_1+1}^n(Y_i - \hat Y_i)^2
\quad (\hat Y_i \text{ calculé avec les paramètres d'entraînement}).$$
Défauts~: (i) partition arbitraire (variabilité)~;
(ii) $\beta$ estimé sur une fraction des données seulement.

\textbf{(2) LOOCV} (exclusion d'une observation)~: pour $i=1,\ldots,n$,
estimer sans l'observation $i$, prédire $\hat y_{i,(-i)}$~:
$$VC_{(n)} = \frac{1}{n}\sum_{i=1}^n(y_i - \hat y_{i,(-i)})^2.$$

\textbf{Formule magique (régression linéaire)}~: un seul ajustement suffit~!
$$VC_{(n)} = \frac{1}{n}\sum_{i=1}^n\left(\frac{y_i - \hat y_i}{1-h_i}\right)^2
\qquad (h_i = \text{levier}).$$

\textbf{(3) Validation croisée à $K$ blocs} ($K=5$ ou $10$)~:
partition $A_1,\ldots,A_K$~; pour chaque $k$, estimer sans $A_k$~:
$$MSE_k = \frac{1}{n_k}\sum_{i\in A_k}\{y_i - \hat y_{i,(-k)}\}^2, \qquad
VC_{(K)} = \frac{1}{K}\sum_{k=1}^K MSE_k.$$
\end{tableau}

\begin{momentR}
\textsf{cv.glm} (librairie \texttt{boot}) sur les 4 modèles de \textsf{Auto}~:
$VC_{(n)}$, $VC_{(5)}$, $VC_{(10)}$ désignent tous M3 (quadratique),
$VC_{(n)}\approx 19.25$, alors qu'une partition unique pouvait suggérer M4.
\end{momentR}

\bloc{3 --- Performance en classification}

\begin{tableau}
Réponse binaire~: le taux d'erreur test devient le
\textbf{taux de mauvaise classification}
$$\frac{1}{n}\sum_{i=1}^n\mathbf{1}_{\{y_i\neq\hat y_i\}},
\qquad \hat y_i = \mathbf{1}_{\{\hat\pi_i > u\}}.$$
Dépend du seuil $u$~! Alternative sans seuil~: l'\textbf{AUC}
(courbe ROC, semaine 5) comme critère de validation croisée.
\end{tableau}

Annoncer~: la validation croisée sera l'outil universel de tout le reste de la
session (choix de $\lambda$, de $k$, de $h$, des paramètres de lissage).

%% ============================================================
%% SEMAINE 10
%% ============================================================
\semaine{10}{Sélection de variables et régularisation~: ridge et lasso}

\noindent\textbf{Objectifs}~: (i) motiver la régularisation par le compromis
biais-variance~; (ii) régression ridge (solution explicite, rétrécissement)~;
(iii) lasso et sélection de variables~; (iv) choix de $\lambda$ par validation croisée.

\bloc{1 --- Le compromis biais-variance}

\begin{tableau}
Moindres carrés~: $\hat\beta$ sans biais, variance minimale
\textbf{parmi les estimateurs sans biais}. Mais le critère pertinent est
l'\textbf{erreur quadratique moyenne}~:
$$E[(\beta_j - \hat\beta_j)^2] = \text{biais}(\hat\beta_j)^2 + Var(\hat\beta_j).$$
Un estimateur \textbf{biaisé} mais de faible variance peut faire mieux~!

\textbf{Régression régularisée}~: minimiser
$$l_\lambda(\beta) = l(\beta) + \lambda\, g(\beta_1,\ldots,\beta_p),
\qquad \lambda\geq0 \text{ (coefficient de pénalité)}.$$
$\lambda\uparrow$ \ $\Rightarrow$ \ biais $\uparrow$, variance $\downarrow$.

Conventions~: $\beta_0$ \textbf{jamais pénalisé}~; variables
\textbf{standardisées} avant l'ajustement (la pénalité n'est pas invariante
d'échelle).
\end{tableau}

\bloc{2 --- Régression ridge}

\begin{tableau}
Pénalité $\ell_2$~: minimiser
$$l^{(R)}_\lambda(\beta) = \sum_{i=1}^n\{Y_i - (\beta_0+\beta_1X_{i1}+\cdots
+\beta_pX_{ip})\}^2 + \lambda\sum_{j=1}^p\beta_j^2.$$
Forme matricielle~: $(Y-X\beta)^\top(Y-X\beta) + \lambda\beta^\top\tilde I_p\beta$
($\tilde I_p$~: identité avec un 0 en position $(1,1)$).

Solution explicite (dérivation identique à la semaine 1)~:
$$\boxed{\hat\beta^{(R)}_\lambda = (X^\top X + \lambda\tilde I_p)^{-1}X^\top Y}$$

Propriétés~:
$$E[\hat\beta^{(R)}_\lambda] = (X^\top X+\lambda\tilde I_p)^{-1}X^\top X\beta
\neq\beta \quad\text{(biaisé si }\lambda>0),$$
$$\lambda = 0~: \hat\beta^{(R)} = \hat\beta~; \qquad
\lambda\to\infty~: \hat\beta_j^{(R)}\to0\ (j\geq1)
\quad\text{= \textbf{rétrécissement} (\textit{shrinkage})}.$$
\end{tableau}

\begin{momentR}
\textsf{Hitters} (\texttt{ISLR}, salaire de joueurs de baseball, 19 variables)~:
\textsf{glmnet(x, y, alpha=0)}. Tracer les trajectoires
$\hat\beta_{j,\lambda}^{(R)}$ en fonction de $\lambda$.
Avec $\lambda=e^6$~: $\sum_j[\hat\beta^{(R)}_{j}]^2 = 5721$ vs $18469$ (MC).
Erreur test~: $124\,716$ (MC) $\to$ $91\,274$ (ridge)~: gain de $\sim27\%$.
\end{momentR}

\bloc{3 --- Régression lasso et choix de $\lambda$}

\begin{tableau}
Pénalité $\ell_1$~: minimiser
$$l^{(L)}_\lambda(\beta) = \sum_{i=1}^n\{Y_i - (\beta_0+\cdots+\beta_pX_{ip})\}^2
+ \lambda\sum_{j=1}^p|\beta_j|.$$
Pas de solution explicite (mais algorithmes efficaces~: \textsf{glmnet},
\textsf{alpha=1}).

\textbf{Propriété clé}~: pour $\lambda$ assez grand, certains
$\hat\beta^{(L)}_{j,\lambda}$ deviennent \textbf{exactement 0}
$\to$ \textbf{sélection de variables} automatique.
(Dessiner l'argument géométrique~: contours elliptiques du SSE et
boule $\ell_1$ en losange vs boule $\ell_2$ ronde.)

\textbf{Choix de $\lambda$}~: validation croisée sur une grille~:
choisir $\lambda$ qui minimise l'erreur test estimée.
En \textrm{R}~: \textsf{cv.glmnet(x, y, alpha=...)}, puis
\textsf{\$lambda.min}.
\end{tableau}

\begin{momentR}
Lasso sur \textsf{Hitters} avec $\lambda = e^2$~: 7 coefficients annulés
(\textsf{HmRun}, \textsf{Runs}, \textsf{RBI}, \textsf{CAtBat}, \textsf{CHits},
\textsf{Assists}, \textsf{NewLeagueN})~: modèle à 12 variables.
\textsf{cv.glmnet}~: $\lambda = e^{2.2}$ (ridge), $\lambda = e^3$ (lasso).
\end{momentR}

%% ============================================================
%% SEMAINE 11
%% ============================================================
\semaine{11}{Régression non-paramétrique~: $k$ plus proches voisins et noyaux}

\noindent\textbf{Objectifs}~: (i) estimateur $k$NN et son compromis biais-variance~;
(ii) distance de Mahalanobis pour plusieurs variables~; (iii) estimateur de
Nadaraya-Watson~; (iv) classification non-paramétrique.

\bloc{1 --- L'estimateur des $k$ plus proches voisins}

Montrer le nuage \textsf{mpg} vs \textsf{weight} (\textsf{Auto})~: relation
clairement non linéaire. On abandonne toute forme fonctionnelle.

\begin{tableau}
Modèle non-paramétrique~:
$$Y_i = f(X_i) + \epsilon_i, \qquad E[\epsilon_i\mid X_i] = 0,
\quad f \text{ inconnue, lisse}.$$
Alors $f(x) = E[Y\mid X=x]$~: la fonction de régression.

\textbf{Estimateur $k$NN}~: moyenne des $Y_i$ des $k$ plus proches voisins~:
$$\hat f(x) = \frac{1}{k}\sum_{i\in N_k(x)}Y_i,$$
$N_k(x)$~: indices des $k$ plus petites distances $d_i(x) = (X_i - x)^2$.

\textbf{Compromis biais-variance} (développement de Taylor autour de $x$)~:
$$\text{biais}[\hat f(x)]\approx\frac{f''(x)}{2k}\sum_{i\in N_k(x)}(E[X_i]-x)^2
\quad(\nearrow \text{ avec } k), \qquad
Var[\hat f(x)] = \frac{\sigma^2}{k} \quad(\searrow \text{ avec } k).$$
Choix de $k$~: minimiser biais$^2$ $+$ variance $\to$ en pratique,
\textbf{validation croisée}.
\end{tableau}

Dessiner deux estimations au tableau~: $k$ petit (courbe erratique),
$k$ grand (courbe trop plate).

\begin{momentR}
\textsf{knn.reg} (librairie \texttt{FNN}) sur \textsf{Auto}~: montrer
$k = 2$, $k = 45$, $k = 300$~; le critère \textsf{PRESS} (LOOCV) donne
$k_{\min} = 45$.
\end{momentR}

\bloc{2 --- Plusieurs variables et estimateur de noyau}

\begin{tableau}
\textbf{Plusieurs variables}~: seule la distance change. La distance
euclidienne échoue si les échelles diffèrent (poids~: 1500--5000~;
cylindrée~: 60--500). \textbf{Distance de Mahalanobis}~:
$$d_i(x) = (X_i - x)^\top V^{-1}(X_i - x), \qquad
V = \text{matrice de var-cov empirique}.$$

\textbf{De $k$NN aux noyaux}~: réécrire $k$NN comme moyenne pondérée~:
$$\hat f(x) = \frac{\sum_i \mathbf{1}_{\{i\in N_k(x)\}}Y_i}
{\sum_i \mathbf{1}_{\{i\in N_k(x)\}}}.$$
Autre voisinage~: $|X_i - x|\leq h$ ($h$ = \textbf{fenêtre})~:
$$\hat f(x) = \frac{\sum_i K\!\left(\frac{X_i-x}{h}\right)Y_i}
{\sum_i K\!\left(\frac{X_i-x}{h}\right)} = \sum_i w_i(x)Y_i,
\qquad K(u) = \mathbf{1}_{\{|u|<1\}} \text{ (noyau uniforme)}.$$
Problème~: poids qui « sautent » $\to$ $\hat f$ discontinue.

\textbf{Nadaraya-Watson}~: noyau continu, p.~ex.\ gaussien
$K(u) = e^{-u^2/2}$ $\to$ poids continus, $\hat f_{NW}$ continue.

\textbf{Choix de $h$}~: LOOCV, avec formule explicite~:
$$cv.n = \frac{1}{n}\sum_{i=1}^n
\left(\frac{Y_i - \hat f(X_i)}{1 - w_i(X_i)}\right)^2.$$
\end{tableau}

\begin{momentR}
\textsf{locpoly} (librairie \texttt{KernSmooth}) sur \textsf{Auto}~:
noyau gaussien, $h_{\min} = 110$ par validation croisée~; comparer
petits et grands $h$.
\end{momentR}

\bloc{3 --- Classification non-paramétrique}

\begin{tableau}
\textsf{Default} (\texttt{ISLR}, $n = 10\,000$)~: $Y=1$ si défaut de carte
de crédit, $X$ = solde. Modèle~: $g(x) = P(Y=1\mid X=x)$, estimé par
$$\hat g(x) = \frac{1}{k}\sum_{i\in N_k(x)}Y_i \in[0,1].$$

Choix de $k$ par validation croisée, deux critères~:
\begin{enumerate}
\item \textbf{taux de mauvaise classification} avec seuil $u$~:
$k = 190$ ($u = 1/2$)~; mais dépend de $u$
($u=1/4$~: $k=290$~; $u=3/4$~: $k=50$)~;
\item \textbf{AUC} (à maximiser, aucun seuil)~: $k = 200$.
\end{enumerate}
\end{tableau}

%% ============================================================
%% SEMAINE 12
%% ============================================================
\semaine{12}{Régression non-linéaire~: polynômes, fonctions par morceaux et splines}

\noindent\textbf{Objectifs}~: (i) cadre des bases de fonctions~; (ii) régression
polynomiale et constantes par morceaux~; (iii) splines~: comptage des degrés de
liberté, base des puissances tronquées, B-splines.

\bloc{1 --- Bases de fonctions et régression polynomiale}

\begin{tableau}
Idée générale~: garder la machinerie \textbf{linéaire en $\beta$} avec des
transformations connues~:
$$Y_i = \beta_0 b_0(X_i) + \beta_1 b_1(X_i) + \cdots + \beta_q b_q(X_i) + \epsilon_i.$$
Tout s'applique~: $\hat\beta = (X^\top X)^{-1}X^\top Y$, IC, tests.

\textbf{Régression polynomiale}~: $b_k(x) = x^k$~:
$$f(x) = \beta_0 + \beta_1 x + \cdots + \beta_p x^p
\qquad (p+1 \text{ \textbf{degrés de liberté}}).$$
Choix de $p$~: AIC, BIC, tests $F$, ou \textbf{validation croisée} (sem.\ 9).

Bandes de confiance~: $\hat f(x) = x_*\hat\beta$, $x_* = (1,x,\ldots,x^p)$~:
$$Var[\hat f(x)] = x_*\,Var[\hat\beta]\,x_*^\top.$$

\textbf{Constantes par morceaux}~: partition
$A = c_0 < c_1 < \cdots < c_K = B$ du domaine,
$b_k(x) = \mathbf{1}_{\{x\in[c_{k-1},c_k]\}}$~:
$$\hat\beta_k = \frac{\sum_i b_k(X_i)Y_i}{n_k} = \text{moyenne locale},
\qquad Var[\hat\beta_k] = \sigma^2/n_k.$$
($X$ discrétisée~: ANOVA à un facteur à $K$ modalités.)
\end{tableau}

\begin{momentR}
\textsf{lm(mpg \textasciitilde{} poly(weight, 2, raw=TRUE))} et
\textsf{lm(mpg \textasciitilde{} cut(weight, 4))} sur \textsf{Auto}~:
superposer polynôme et fonction en escalier sur le nuage de points.
\end{momentR}

\bloc{2 --- Splines~: définition et degrés de liberté}

\begin{tableau}
\textbf{Splines}~: polynômes par morceaux avec contraintes de régularité aux
\textbf{n\oe uds} (\textit{knots}).

Cas simple ($K=2$ intervalles, degré 3)~:
$$f(x) = \begin{cases}
\beta_0 + \beta_1x + \beta_2x^2 + \beta_3x^3 & x\in[c_0,c_1)\\
\gamma_0 + \gamma_1x + \gamma_2x^2 + \gamma_3x^3 & x\in[c_1,c_2)
\end{cases}$$
$4+4 = 8$ paramètres~; imposer $f$, $f'$, $f''$ continues en $c_1$
($3$ contraintes) $\Rightarrow$ $8 - 3 = 5$ degrés de liberté.

\textbf{Cas général}~: $K$ intervalles, degré $p$~:
$$\underbrace{K(p+1)}_{\text{paramètres}} - \underbrace{p(K-1)}_{\text{contraintes}}
= K + p \ \text{ degrés de liberté}.$$

\textbf{Base des puissances tronquées} ($K+p$ fonctions)~:
$$1,\ x,\ \ldots,\ x^p,\quad (x-c_k)_+^p \ (k = 1,\ldots,K-1),$$
où $(x-c)_+^p = (x-c)^p\,\mathbf{1}_{\{x\geq c\}}$.
Vérifier au tableau que $(x-c)_+^p$ est bien $p-1$ fois continûment dérivable
en $c$.
\end{tableau}

\bloc{3 --- B-splines}

\begin{tableau}
Puissances tronquées~: simples mais \textbf{numériquement instables}
($X^\top X$ mal conditionnée). Les logiciels utilisent les \textbf{B-splines}~:
base équivalente, à \textbf{support compact} (chaque fonction non nulle sur
$\leq p+1$ intervalles).

Récurrence de Cox-de Boor (sur les n\oe uds augmentés $\{t_j\}$)~:
$$B_{j,0}(x) = \mathbf{1}_{\{t_j\leq x < t_{j+1}\}},$$
$$B_{j,d}(x) = \frac{x-t_j}{t_{j+d}-t_j}B_{j,d-1}(x)
+ \frac{t_{j+d+1}-x}{t_{j+d+1}-t_{j+1}}B_{j+1,d-1}(x).$$
Propriétés~: $B_{j,p}\geq0$~; $\sum_j B_{j,p}(x) = 1$ (partition de l'unité)~;
matrice de plan creuse par bandes $\Rightarrow$ calculs stables et rapides.
\end{tableau}

\begin{momentR}
\textsf{bs} (librairie \texttt{splines})~:
\textsf{lm(mpg \textasciitilde{} bs(weight, knots=c(2200,2800,3600)), data=Auto)}.
Tracer les fonctions de base B-splines elles-mêmes, puis l'ajustement.
\end{momentR}

%% ============================================================
%% SEMAINE 13
%% ============================================================
\semaine{13}{Splines naturelles, splines de lissage et modèles additifs}

\noindent\textbf{Objectifs}~: (i) splines naturelles~; (ii) splines de lissage
(pénalisées) et lien avec ridge~; (iii) degrés de liberté effectifs~;
(iv) le modèle additif gaussien~: définition et estimation pénalisée.

\bloc{1 --- Splines naturelles et splines de lissage}

\begin{tableau}
\textbf{Splines naturelles}~: comportement erratique des splines aux
extrémités $\to$ imposer $f$ \textbf{linéaire} sur $[c_0,c_1)$ et
$[c_{K-1},c_K)$. Contraintes supplémentaires $=$ degrés de liberté libérés
(réinvestis en n\oe uds intérieurs). En \textrm{R}~: \textsf{ns(weight, df=4)}.

\textbf{Splines de lissage}~: contre le surajustement, pénaliser la rugosité~:
$$\min_\beta\ \sum_{i=1}^n\Bigl\{Y_i - \sum_k\beta_kb_k(X_i)\Bigr\}^2
+ \lambda\int_A^B f''(t)^2\,dt.$$
Comme $f = \sum_k\beta_kb_k$~: pénalité $= \lambda\beta^\top M\beta$,
$M_{lq} = \int b''_l(t)b''_q(t)\,dt$~:
$$\boxed{\hat\beta_\lambda = (X^\top X + \lambda M)^{-1}X^\top Y}$$
\textbf{Analogie parfaite avec ridge~!} (sem.\ 10)
\begin{itemize}
\item $\lambda = 0$~: pas de pénalité, surajustement possible~;
\item $\lambda\to\infty$~: $f''\equiv0$~: droite des moindres carrés.
\end{itemize}

\textbf{Degrés de liberté effectifs}~: $\hat Y = A_\lambda Y$,
$A_\lambda = X(X^\top X+\lambda M)^{-1}X^\top$~:
$$dle = \mathrm{Tr}(A_\lambda) \in\ ]2,\ K+p], \quad
\text{décroît quand } \lambda \text{ croît}.$$
\end{tableau}

\begin{momentR}
\textsf{smooth.spline(Auto\$weight, Auto\$mpg, cv=TRUE)}~:
$\lambda$ choisi par validation croisée, $dle\approx10.4$.
Faire varier \textsf{df} pour montrer l'effet du lissage.
\end{momentR}

\bloc{2 --- Le modèle additif}

\begin{tableau}
Problème en dimension $p$ grande~: \textbf{fléau de la dimension}
(voisinages vides, explosion du nombre de fonctions de base).

\textbf{Modèle additif}~: compromis flexibilité/interprétabilité~:
$$Y_i = \alpha + f_1(X_{i1}) + f_2(X_{i2}) + \cdots + f_p(X_{ip}) + \epsilon_i,$$
chaque $f_j$ lisse et inconnue, représentée par une base de splines~:
$f_j(x) = \sum_{k=1}^{K_j}b_{jk}(x)\beta_{jk}$.

\textbf{Identifiabilité}~: une constante peut circuler entre $\alpha$ et les
$f_j$ $\to$ contrainte de centrage
$$\sum_{i=1}^n f_j(X_{ij}) = 0, \qquad j=1,\ldots,p.$$
\end{tableau}

\bloc{3 --- Estimation pénalisée et inférence}

\begin{tableau}
Bases riches $+$ pénalités~: minimiser
$$\|Y - X\beta\|^2 + \sum_{j=1}^p\lambda_j\,\beta^\top S_j\,\beta
\qquad (\beta^\top S_j\beta = \textstyle\int f_j''(t)^2dt),$$
$$\hat\beta = \Bigl(X^\top X + \sum_j\lambda_jS_j\Bigr)^{-1}X^\top Y.$$
$\lambda_j\to\infty$~: $f_j$ linéaire~; $\lambda_j = 0$~: $f_j$ libre.

$\hat Y = AY$~: degrés de liberté effectifs $\mathrm{Tr}(A)$,
\textbf{décomposables par fonction} $f_j$.

\textbf{Choix des $\lambda_j$}~:
$$OCV(\lambda) = \frac1n\sum_i\left(\frac{y_i-\hat f_\lambda(x_i)}{1-A_{ii}}\right)^2,
\qquad
GCV(\lambda) = \frac{n\sum_i(y_i-\hat f_\lambda(x_i))^2}{(n-\mathrm{Tr}(A))^2}.$$

\textbf{Inférence}~: forme ridge $\Rightarrow$ $\hat\beta$ gaussien~;
approche bayésienne~: $V_\beta = \sigma^2(X^\top X+\sum_j\lambda_jS_j)^{-1}$,
intervalles $\hat f_j(x)\pm z_{\alpha/2}\sqrt{\widehat{Var}[\hat f_j(x)]}$
(bonne couverture fréquentiste). Test de $H_0: f_j = 0$~: Wald
($\chi^2$ si $\phi$ connu, $F$ sinon).
\end{tableau}

\begin{momentR}
Premier modèle additif~: \textsf{gam(mpg \textasciitilde{} s(weight) +
s(horsepower), data=Auto)} (librairie \texttt{mgcv})~;
\textsf{plot(mod)} pour les courbes $\hat f_j$ avec bandes de confiance.
\end{momentR}

%% ============================================================
%% SEMAINE 14
%% ============================================================
\semaine{14}{Modèles additifs généralisés (GAM) et synthèse du cours}

\noindent\textbf{Objectifs}~: (i) définir le GAM (structure additive $+$ famille
exponentielle)~; (ii) algorithme d'ajustement~; (iii) UBRE/GCV et comparaison
de modèles~; (iv) pratique avec \texttt{mgcv}~; (v) synthèse générale du cours.

\bloc{1 --- Définition et ajustement des GAM}

\begin{tableau}
\textbf{GAM} $=$ MLG (sem.\ 7) $+$ structure additive (sem.\ 13)~:
$Y_i\sim$ famille exponentielle, $E[Y_i] = \mu_i$,
$$g(\mu_i) = \alpha + f_1(X_{i1}) + \cdots + f_p(X_{ip}).$$
Cas particuliers~: GAM logistique (binaire), GAM Poisson (dénombrements,
offset au besoin).

\textbf{Estimation}~: déviance pénalisée~:
$$\hat\beta = \arg\min_\beta\Bigl\{D(y,\mu)
+ \sum_{j=1}^p\lambda_j\,\beta^\top S_j\,\beta\Bigr\}.$$

\textbf{Algorithme} (moindres carrés pondérés itératifs pénalisés)~;
à l'itération $m$~:
\begin{enumerate}
\item variable ajustée~:
$z_i^{(m)} = \eta_i^{(m)} + (y_i-\mu_i^{(m)})
\left(\frac{\partial\eta_i}{\partial\mu_i}\right)^{(m)}$~;
\item poids~: $w_i^{(m)} = \frac{1}{V(\mu_i^{(m)})}
\left(\frac{\partial\mu_i}{\partial\eta_i}\right)^2_{(m)}$~;
\item ajuster un modèle additif pénalisé pondéré à $z^{(m)}$~;
\item répéter jusqu'à convergence.
\end{enumerate}
\end{tableau}

\bloc{2 --- Choix des paramètres de lissage et pratique avec R}

\begin{tableau}
\textbf{$\phi$ connu} (logistique, Poisson)~: critère \textbf{UBRE}~:
$$UBRE(\lambda) = \frac{D(y,\hat\mu)}{n} + \frac{2\phi\,\mathrm{Tr}(A)}{n}~;$$
\textbf{$\phi$ inconnu}~: \textbf{GCV}~:
$$GCV(\lambda) = \frac{n\,D(y,\hat\mu)}{(n - \mathrm{Tr}(A))^2}.$$

\textbf{Comparaison de modèles}~:
$AIC = -2\,l(\hat\beta;y) + 2\,\mathrm{Tr}(A)$
(dle au lieu du nombre de paramètres), GCV, proportion de déviance expliquée
$(D_{null}-D)/D_{null}$.

\textbf{Diagnostic}~: $dle(f_j)\approx1$ $\Rightarrow$ $f_j$ remplaçable par
un terme linéaire~: le GAM guide la spécification d'un MLG paramétrique.
\end{tableau}

\begin{momentR}
Reprendre les jeux de données de toute la session avec \texttt{mgcv}~:
\begin{itemize}
\item \textsf{gam(mpg \textasciitilde{} s(weight) + s(horsepower), data=Auto)}~;
\item \textsf{gam(y \textasciitilde{} s(SOI), family=binomial, data=quilpie)}~;
\item \textsf{gam(Cases \textasciitilde{} s(AgeNum) + City + offset(log(Pop)),
family=poisson, data=danishlc)}.
\end{itemize}
\textsf{summary}~: dle et test de nullité par fonction~;
\textsf{plot}~: courbes avec bandes de confiance.
\end{momentR}

\bloc{3 --- Synthèse du cours}

Construire le tableau récapitulatif \textbf{avec} les étudiants~: donner les
lignes/colonnes et leur demander de remplir chaque case.

\begin{tableau}
\begin{tabular}{lll}
\toprule
 & \textbf{Relation linéaire} & \textbf{Relation flexible} \\
\midrule
Réponse normale & régression linéaire & polynômes, splines, $k$NN,\\
 & (+ ridge/lasso) & Nadaraya-Watson, modèle additif \\
Réponse binaire & régression logistique & GAM logistique, $k$NN \\
Dénombrement & régression Poisson & GAM Poisson \\
Famille exponentielle & MLG & GAM \\
\bottomrule
\end{tabular}

\medskip
\textbf{Fil conducteur}~: l'estimation de l'erreur de prévision
(\textbf{validation croisée}) pour choisir la complexité~:
degré $p$, nombre de voisins $k$, fenêtre $h$, pénalités $\lambda$, $\lambda_j$.
\end{tableau}

Terminer par une période de questions et de révision en vue de l'examen final.

\end{document}
